\section{Proposal：危険度連動eHMI}
\label{sec:proposal}

\subsection{処理の流れ}

各AGVについて，毎フレーム以下を実行する．

\begin{enumerate}[leftmargin=2em]
  \item 参加者の視線方向から\textbf{動的基準線}（太線ゾーン）を更新
  \item 人へのPTTC $T_p$，経路--基準線TTC $T_c$，水平距離 $d$ を算出
  \item 重み付き和で危険度スコア $R \in [0,1]$ を計算（表示は $d\le D_\text{max}$ のみ）
  \item $R$ を色相・彩度・明度・不透明度に変換し，経路（または停止線）を描画
\end{enumerate}

\subsection{動的基準線（視線方向）}

目的地最短経路ではなく，\textbf{現在見ている方向}に注目領域を置く．
HMDのforwardをXZ平面に投影した $\mathbf{f}$ からレイキャストし，
障害物との交点 $\mathbf{h}$ を求める．床面上の基準軸は
\begin{equation}
  \mathbf{a}_\text{start} = \mathbf{p}, \quad
  \mathbf{a}_\text{end} = \mathrm{flatten}(\mathbf{h})
\end{equation}
半幅 $w=\SI{0.5}{m}$（全幅\SI{1.0}{m}）の帯状ゾーンとして，
AGV経路との交差判定に用いる．
立ち止まって別方向を見た場合も，視線に追従して危険評価が更新される．

\subsection{危険度スコア $R$}

\paragraph{表示可否}
表示ゲートは水平距離のみとする（$T_\text{max}$ はスコア正規化専用）．
\begin{equation}
  \text{visible} = (d \le D_\text{max})
\end{equation}
$T_\text{max}=\SI{4.0}{s}$，$D_\text{max}=\SI{13.6}{m}=(v_\text{AGV,max}+v_\text{ped})\times T_\text{max}$．

\paragraph{三要因}
人へのPTTC $T_p=d/v_\text{close}$（$v_\text{close}=\max(0,\hat{\mathbf{r}}\cdot(\mathbf{v}_\text{AGV}-\mathbf{v}_\text{ped}))$；近づかないとき $\infty$），
経路--基準線TTC $T_c=s/\max(v,0.1)$，水平距離 $d$．

\paragraph{統合（重み付き和）}
\begin{align}
  R_p &= f(T_p),\quad R_c=f(T_c),\quad
  R_d=(1-d/D_\text{max})^{\gamma}\ \text{（$d\ge D_\text{max}$ なら $0$）} \\[2pt]
  f(T) &= (1-T/T_\text{max})^{\gamma}\ \text{（$T=\infty$ または $T>T_\text{max}$ なら $0$）} \\[2pt]
  R &= \min(1,\, w_p R_p + w_c R_c + w_d R_d)
\end{align}
$\gamma=0.6$，$w_p=0.55$，$w_c=0.30$，$w_d=0.15$。PTTC は相対速度 $\mathbf{v}_\text{AGV}-\mathbf{v}_\text{ped}$。

\subsection{視覚マッピング（Proposedのみ）}

色相・不透明度はともに $R$ に対して連続変化する．
\begin{align}
  H(R) &= (1-R) \times 180° \quad \text{（青緑$\rightarrow$赤）} \\
  S(R) &= 0.4 + 0.6R, \quad V(R) = 0.5 + 0.3R \\
  \alpha(R) &= 0.35 + 0.65R \quad \text{（半透明$\rightarrow$不透明）}
\end{align}

危険度が上がると \textbf{青緑・半透明 $\rightarrow$ 赤・不透明} へ連続的に変化し，
「だんだん危なくなっている」ことが視覚的に伝わる．

\subsection{経路描画}

\begin{itemize}[leftmargin=2em]
  \item \textbf{走行中}: 残存経路をテーパー付きライン（根元\SI{0.65}{m}，先端\SI{0.12}{m}）で表示
  \item \textbf{停止中}: 停止線（全幅\SI{1.0}{m}）を表示
  \item \textbf{表示距離}: 被験者から水平\SI{13.6}{m}以内のAGVのみ
  \item \textbf{描画優先}: $\mathrm{sortingOrder}=\mathrm{round}(R\times100)$（危険なAGVを手前に）
\end{itemize}

\begin{table}[htbp]
  \centering
  \caption{3条件の対照（Proposal節のまとめ）}
  \label{tab:conditions}
  \begin{tabular}{lccc}
    \toprule
    項目 & Baseline & No-AR & Proposed \\
    \midrule
    経路表示 & あり & なし & あり \\
    危険度計算 & なし（$R=1$固定） & --- & PTTC + 経路TTC + 近接 \\
    色 & 固定青 & --- & HSV連動 \\
    不透明度 & 1.0固定 & --- & 連続（$0.35$–$1.0$） \\
    \bottomrule
  \end{tabular}
\end{table}

\subsection{Proposalとしての妥当性}

\paragraph{再現可能性}
$R$ は PTTC，経路TTC，距離，速度から\textbf{決定論的}に算出される．
被験者や実験者の主観が刺激生成に介入しない．

\paragraph{Baseline / No-AR との分離}
Proposedだけが $R$ に連動するため，
\textbf{危険度エンコーディングの効果}を他条件と比較して評価できる．

\paragraph{パラメータの根拠}
$D_\text{max}$ は表示ゲートと近接スコアの共通上限で，
$(v_\text{AGV,max}+v_\text{ped})\times T_\text{max}$ から導出する．
$T_\text{max}$ は時間スコアの正規化専用である．
